\chapter{Notation and Conventions}
\label{app:notation-conventions}

This appendix collects all notation, conventions, and terminological
distinctions used across Volume~III into a single reference.  Where a
convention is shared with Volumes~I--II, we state it here for
self-containedness and cite the canonical source.

%% =====================================================================
\section{Grading and sign conventions}
\label{sec:nc-grading}

\begin{itemize}[nosep]
\item All grading is \textbf{cohomological}: differentials have degree $+1$.
\item The bar construction uses desuspension $s^{-1}$
      (matching Volumes~I--II).
\item Signs follow the Koszul rule throughout.  For the curved
      $\Ainf$-structure: $\mu_1^2(a) = [\mu_0, a]$ (commutator, minus sign).
\item Tensor products of dg categories are derived unless stated otherwise.
\item ``Equivalence'' of dg categories means quasi-equivalence
      (equivalence of associated $\infty$-categories).
\end{itemize}

%% =====================================================================
\section{The \texorpdfstring{$\kappa$}{kappa}-spectrum}
\label{sec:nc-kappa}

A single CY manifold $X$ admits \emph{multiple} chiral algebras, each
arising from a different mathematical construction ($\Phi$
produces exactly one; the others come from independent constructions).
Each construction has its own characteristic exponent.
Bare ``$\kappa$'' without subscript is \textbf{forbidden} in Vol~III.
The four subscripted invariants fall into two types:

\medskip\noindent
\textbf{Manifold invariants} (topological, fixed by the geometry of $X$,
independent of algebraization):
\begin{center}
\begin{tabular}{@{}lp{4.5cm}p{5cm}l@{}}
\toprule
\textbf{Symbol} & \textbf{Name} & \textbf{Definition / Source} & \textbf{K3$\times E$} \\
\midrule
$\kappa_{\mathrm{cat}}$ & Categorical &
  Holomorphic Euler characteristic $\chi(\cO_X)$ & 2 \\
$\kappa_{\mathrm{fiber}}$ & Fiber / lattice &
  Lattice rank or fiber-structure invariant
  (Mukai lattice rank for K3) & 24 \\
\bottomrule
\end{tabular}
\end{center}

\noindent
\textbf{Algebraization invariants} (depend on which chiral/BKM algebra
is constructed from~$X$):
\begin{center}
\begin{tabular}{@{}lp{4.5cm}p{5cm}l@{}}
\toprule
\textbf{Symbol} & \textbf{Name} & \textbf{Definition / Source} & \textbf{K3$\times E$} \\
\midrule
$\kappa_{\mathrm{ch}}$ & Chiral &
  From the chiral algebra $A_\cC = \Phi(\cC)$; the weight
  appearing in the modular trace $\Tr_{A_\cC}$ & 3 \\
$\kappa_{\mathrm{BKM}}$ & Borcherds--Kac--Moody &
  Weight of the BKM automorphic form
  ($\Delta_5$ for K3$\times E$) & 5 \\
\bottomrule
\end{tabular}
\end{center}

\noindent
Stating that ``algebraizations share $\kappa_{\mathrm{cat}}$'' is vacuous:
$\kappa_{\mathrm{cat}}$ and $\kappa_{\mathrm{fiber}}$ are topological
invariants of the manifold and cannot vary between algebraizations.

\medskip\noindent
\textbf{Forbidden subscripts}: \texttt{global}, \texttt{BPS},
\texttt{eff}, \texttt{total}, \texttt{naive}, \texttt{MacMahon}.
If the subscript is an indexing variable over the approved set,
write $\kappa_\bullet$.

\medskip\noindent
\textbf{Key relations}:
\begin{itemize}[nosep]
\item Koszul conductor: $\kappa_{\mathrm{ch}}(A) +
      \kappa_{\mathrm{ch}}(A^!) = K$ (family-dependent; $K = 0$ for
      K3 free-field/Kac--Moody branch).
\item Flops preserve $\kappa_{\mathrm{ch}}$:
      $\kappa_{\mathrm{ch}}(A_X) = \kappa_{\mathrm{ch}}(A_{X^+})$
.
\item $\kappa_{\mathrm{ch}} = \chi^{\CY}$ at $d = 2$ (proved via
      Serre duality $S_\cC = [2]$); conjectured at $d \ge 3$.
\end{itemize}

%% =====================================================================
\section{The \texorpdfstring{$\En$}{E\_n} hierarchy}
\label{sec:nc-en}

\begin{center}
\begin{tabular}{@{}llll@{}}
\toprule
\textbf{Level} & \textbf{Operad} & \textbf{Structure} & \textbf{CY source} \\
\midrule
$\Eone$ & Little 1-disks & Associative / ordered factorization
  on $\C \times \R$ & Native for $d \ge 3$ (hol.\ CS) \\
$\Etwo$ & Little 2-disks & Braided factorization on
  $\C \times \C$ & $d = 2$ ($S^2$-framing); $d = 3$ via
  Drinfeld center \\
$E_3$ & Little 3-disks & Framed 3-disk factorization &
  6d hol.\ theory on $\C^3$ \\
$\Einf$ & $\mathrm{Comm}$ & Symmetric / commutative &
  Kills Hopf / quantum group structure \\
\bottomrule
\end{tabular}
\end{center}

\noindent
\textbf{Key relationships}:
\begin{itemize}[nosep]
\item Dunn additivity: $\Eone \otimes \Eone \simeq \Etwo$, \;
      $\Eone \otimes \Etwo \simeq E_3$.
\item $\Eone \to \Etwo$ promotion:
      $Z(\Rep^{\Eone}(A)) = \Rep^{\Etwo}(\Zder(A))$ (Drinfeld center).
\item $\Etwo \to E_3$ promotion: derived center (higher Deligne),
      \emph{not} iterated Drinfeld center.
\item $\Etwo$ braiding is \emph{not} symmetric: $\Etwo \neq \Einf$.
\item $\Einf$ averaging $\mathrm{av}\colon B^{\mathrm{ord}} \to B^{\Einf}$
 kills the $R$-matrix: $\mathrm{av}(r(z)) = \kappa_{\mathrm{ch}}$.
\item $\Eone$-stabilization: the $\En$-cascade $\Eone \to \Etwo \to E_3$
      terminates at $n = 3$ for CY inputs; $E_n \simeq E_3$ for $n \ge 3$
      at the chain level.
\end{itemize}

\noindent
\textbf{Dimension--$\En$ correspondence}:
\begin{center}
\begin{tabular}{@{}cll@{}}
\toprule
$d$ & Operadic level & Example \\
\midrule
1 & $\Einf$ (commutative) & Proved, trivial \\
2 & $\Etwo$ (braided) & K3: CY-A$_2$ proved \\
$\ge 3$ & $\Eone$ (ordered) & K3$\times E$: CY-A$_3$ programme \\
\bottomrule
\end{tabular}
\end{center}

%% =====================================================================
\section{Shadow classes G/L/C/M}
\label{sec:nc-shadow}

The shadow class of a chiral algebra $A$ is determined by the
\emph{full} shadow obstruction tower $\{S_k\}_{k \ge 2}$, never by
generator counting or non-formality alone.

\begin{center}
\begin{tabular}{@{}lllll@{}}
\toprule
\textbf{Class} & \textbf{Shadow depth} & \textbf{$\Ainf$ corrections}
  & \textbf{QGL analytic type} & \textbf{Example} \\
\midrule
G (Gaussian) & 0 & None ($m_k = 0$, $k \ge 3$) & Polynomial & Free field \\
L (Lie) & Finite & Finite truncation & Rational & Kac--Moody \\
C (Calabi) & Finite & Charge conservation kills $d_4$ & Convergent
  & Fermionic $bc$ \\
M (Mock) & $\infty$ & All $\delta^{(k)} \neq 0$ & Gevrey-1 (Borel)
  & Virasoro, $\cW_{1+\infty}$ \\
\bottomrule
\end{tabular}
\end{center}

\noindent
The shadow invariants $S_k$ are the $\Ainf$ coproduct correction
coefficients:
\[
  \Delta^{\Ainf} = \Delta^{\mathrm{Yangian}} + \hbar^2 \delta^{(3)}
  + \hbar^3 \delta^{(4)} + \cdots,
\]
with $\delta^{(k)}$ having coefficient $S_k$.

\medskip\noindent
$E_3$ bar cohomology:
$\dim H^*(B_{E_3}(A)) = (1+t)^{3g}$ for classes L and C;
\textbf{fails} for class M ($d_4$ survives, infinite-dimensional;
). Chain-level dimension for all classes: $P(q)^{3g}$.

%% =====================================================================
\section{Three centers}
\label{sec:nc-centers}

Three distinct notions of ``center'' appear in this volume.
They must \emph{never} be conflated.

\begin{center}
\begin{tabular}{@{}p{3.5cm}p{5cm}l@{}}
\toprule
\textbf{Object} & \textbf{Definition} & \textbf{Level} \\
\midrule
Drinfeld center $Z(\cC)$ & Category of half-braidings on the
  monoidal category $\cC$; objects are pairs
  $(X, \sigma_{X,-})$ & Categorical \\[4pt]
Derived center $\Zder(A)$ & Hochschild cochain complex
  $\mathrm{CH}^*(A,A)$; the ``bulk algebra''
  of chiral CFT & Algebraic \\[4pt]
M\"{u}ger center $Z_2(\cC)$ & Full subcategory of objects
  $X \in \cC$ with trivial double braiding
  $c_{Y,X} \circ c_{X,Y} = \id_{X \otimes Y}$
  for all $Y$ & Categorical (degeneracy locus) \\
\bottomrule
\end{tabular}
\end{center}

\noindent
\textbf{Key relationships}:
\begin{itemize}[nosep]
\item $Z(\Rep^{\Eone}(A)) \simeq \Rep^{\Etwo}(\Zder(A))$
      (Drinfeld center categorifies derived center).
\item M\"{u}ger center $Z_2(\cC) = \Vect$ iff $\cC$ is
      \emph{non-degenerate} (modular).
\item For K3 Heisenberg: Drinfeld center is 49-dimensional
      ($24 + 24 + 1$); $>92\%$ invisible to local charts
      (center-hocolim obstruction).
\end{itemize}

%% =====================================================================
\section{Spectral vs.\ worldsheet parameters}
\label{sec:nc-z}

Two mathematically distinct parameters are both conventionally
denoted $z$.  Conflation is the source of the adversarial
``$z = 0$ singularity'' objection.

\begin{center}
\begin{tabular}{@{}lp{5cm}l@{}}
\toprule
\textbf{Parameter} & \textbf{Context} & \textbf{Behaviour at 0} \\
\midrule
$z_{\mathrm{spec}}$ & Yangian spectral parameter: shift
  $u \to u - z$ in the transfer matrix &
  $\Delta_0$ cocommutative; admits antipode \\[4pt]
$z_{\mathrm{ws}}$ & Worldsheet / OPE insertion coordinate:
  $T(z)T(w) \sim c/2 \cdot (z - w)^{-4} + \cdots$ &
  Poles from OPE singularities \\
\bottomrule
\end{tabular}
\end{center}

\noindent
Setting $z = 0$ in the Drinfeld coproduct $\Delta_z$ removes the
spectral shift (no OPE singularity).  Setting $z \to w$ in the OPE
$T(z)T(w)$ produces poles.  Before any $z = 0$ argument, one must
state whether $z$ is spectral or worldsheet.

%% =====================================================================
\section{Ordering conventions}
\label{sec:nc-ordering}

Bare ``ordered'' is \textbf{forbidden}.  Three distinct
orderings appear:

\begin{center}
\begin{tabular}{@{}llll@{}}
\toprule
\textbf{Term} & \textbf{Type} & \textbf{Operation}
  & \textbf{Algebraic structure} \\
\midrule
Labeled-ordered & Combinatorial & $\Eone$ bar
  $B^{\mathrm{ord}}$ (deconcatenation) & Hopf \\
Time-ordered & Analytic / radial & OPE on $\C$ with
  radial ordering & Vertex algebra \\
Normally-ordered & Wick & $\mathopen{:}\phi(z)\phi(w)\mathclose{:}$
  (creation left, annihilation right) & Free field \\
\bottomrule
\end{tabular}
\end{center}

%% =====================================================================
\section{Bar complex and Koszul duality}
\label{sec:nc-bar}

\begin{center}
\begin{tabular}{@{}lp{8cm}@{}}
\toprule
\textbf{Symbol} & \textbf{Meaning} \\
\midrule
$B(A)$ & Bar complex of $A$: a factorization \emph{coalgebra} on $\Ran(X)$ \\
$B^{\mathrm{ord}}(A)$ & $\Eone$ ordered bar complex
  $T^c(s^{-1}\bar{A})$: deconcatenation coproduct,
  habitat of the Hopf structure \\
$B_{\En}(A)$ & $\En$-bar complex: the bar construction with
  respect to the $\En$-operad.  $B_{\Eone} = B^{\mathrm{ord}}$.
  $B_{E_3}$ gives cohomology $(1+t)^{3g}$ for classes L, C \\
$\Omega(C)$ & Cobar construction: $\Omega(B(A)) \simeq A$
  (bar-cobar inversion, Vol~I Theorem~B).
  This is \emph{not} Koszul duality \\
$A^! = (H^*(B(A)))^\vee$ & Koszul dual algebra (linear dual of bar cohomology) \\
$D_{\Ran}(B(A))$ & Verdier dual: $D_{\Ran}(B(A)) \simeq B(A^!)$
  (Vol~I Theorem~A) \\
$\Theta_A$ & Chiral characteristic form / modular trace datum \\
$r(z) = \Res^{\mathrm{coll}}_{0,2}(\Theta_A)$ & Collision
  $r$-matrix: pole orders one less than OPE (bar extracts
  $d\log(z_i - z_j)$) \\
\bottomrule
\end{tabular}
\end{center}

%% =====================================================================
\section{CY categories and dimension}
\label{sec:nc-cy}

\begin{center}
\begin{tabular}{@{}lp{8cm}@{}}
\toprule
\textbf{Symbol} & \textbf{Meaning} \\
\midrule
$\CY_d\text{-}\Cat$ & $\infty$-category of smooth
  $d$-dimensional Calabi--Yau categories \\
$d$ & CY dimension $= \dim_\C X$ for a smooth CY manifold $X$.
 \emph{Not} the real dimension $2d$ \\
$n$ & Complex dimension of the ambient space, when $\MF(W)$ is
  involved: $W \colon \bA^n \to \bA^1$ gives $\CY_{n-2}$
 \\
$D^b(\Coh(X))$ & Bounded derived category of coherent sheaves on $X$ \\
$\Fuk(X)$ & Fukaya category of $X$ \\
$\MF(W)$ & Matrix factorization category of
  $W \colon \bA^n \to \bA^1$; CY dimension $= n - 2$ \\
$\HH_*(A)$ & Hochschild homology of $A$ \\
$\HC^-_d(\cC)$ & Negative cyclic homology: the CY trace target
 (not bare $\HH_d \to k$) \\
$\Phi \colon \CY_d\text{-}\Cat \to \En\text{-}\mathrm{ChirAlg}$
  & CY-to-chiral functor (Theorem CY-A) \\
$A_\cC = \Phi(\cC)$ & Chiral algebra of CY category $\cC$ \\
\bottomrule
\end{tabular}
\end{center}

\noindent
\textbf{MF dimension table}:
\begin{center}
\begin{tabular}{@{}cll@{}}
\toprule
$n$ (variables) & $\CY_{n-2}$ & Example \\
\midrule
2 & CY$_0$ (semisimple) & ADE surface singularities \\
3 & CY$_1$ & Elliptic curves \\
4 & CY$_2$ & Quartic K3 \\
5 & CY$_3$ & Fermat quintic \\
\bottomrule
\end{tabular}
\end{center}

%% =====================================================================
\section{Structure functions}
\label{sec:nc-structure}

Two conventions for structure functions coexist in the
quantum group literature:

\begin{center}
\begin{tabular}{@{}llll@{}}
\toprule
\textbf{Symbol} & \textbf{Parameter} & \textbf{Type}
  & \textbf{Context} \\
\midrule
$g_{ij}(u)$ & Additive spectral $u$ & Rational &
  Yangian / affine; poles at $u = \pm h_i$ \\
$G_{ij}(x)$ & Multiplicative $x = e^u$ & Trigonometric &
  Quantum group / toroidal; poles at $x = q^{\pm h_i}$ \\
\bottomrule
\end{tabular}
\end{center}

\noindent
The rational structure function for charge-$(i,j)$ is
$g_{ij}(u) = \prod_\alpha (u - \alpha_{ij}) / \prod_\beta (u - \beta_{ij})$,
where the numerator zeros and denominator poles are determined by
the Serre relations.

\medskip\noindent
\textbf{Unitarity}: $g(u) \cdot g(-u) = 1$ (algebraic, no reality
assumption; follows from Mukai signature encoded in $\omega^{ij}$).

\medskip\noindent
\textbf{Koszul inversion}: $G(x; q^{-1}) = 1/G(x; q)$.

%% =====================================================================
\section{Quantum groups and \texorpdfstring{$R$}{R}-matrices}
\label{sec:nc-qg}

\begin{center}
\begin{tabular}{@{}lp{8cm}@{}}
\toprule
\textbf{Symbol} & \textbf{Meaning} \\
\midrule
$\Uq(\frakg)$ & Drinfeld--Jimbo quantum group \\
$\Uhbar(\frakg)$ & $\hbar$-adic deformation ($\hbar = \log q$) \\
$\Yg$ & Yangian of $\frakg$ \\
$Y(\frakg_{K3})$ & K3 Yangian: rank 24, Mukai signature $(4,20)$,
  degree-$(24,24)$ structure function \\
$U_{q,t}(\widehat{\widehat{\fgl}}_1)$ & Quantum toroidal
  $\fgl_1$ (Miki automorphism acts here, not on general
 $E_3$ algebras) \\
$R(z)$ & Yang $R$-matrix (spectral parameter $z_{\mathrm{spec}}$) \\
$R_{\mathrm{ch}}(u,v)$ & Two-parameter chiral $R$-matrix:
  $R_{\mathrm{ch}} = R_1(u) R_2(v) R_{12}(u-v)$
  (Zamolodchikov factorization) \\
$\Delta_z$ & Drinfeld coproduct (spectral parameter shift) \\
$S$ & Antipode: $S(T_n) = (2\Psi - 3)T_n$ at spin~2 \\
$\varepsilon$ & Counit (vacuum projection) \\
$K(z)$ & $K$-matrix: $K(z) = e^{-z\, d/du}$ (boundary reflection) \\
$\CoHA(\cC)$ & Cohomological Hall algebra: associative, \emph{not}
 a chiral algebra \\
\bottomrule
\end{tabular}
\end{center}

%% =====================================================================
\section{Factorization and operadic notation}
\label{sec:nc-fact}

\begin{center}
\begin{tabular}{@{}lp{8cm}@{}}
\toprule
\textbf{Symbol} & \textbf{Meaning} \\
\midrule
$\Ran(X)$ & Ran space of non-empty finite subsets of $X$ \\
$\Ran^{\mathrm{ord}}(X)$ & Ordered Ran space (sequences, not sets) \\
$\Fact(\cC)$ & Factorization category \\
$\FM_k(\C)$ & Fulton--MacPherson compactification of
  $\mathrm{Conf}_k(\C)$; blowup along diagonals,
  \emph{not} the complement \\
$\SC^{\mathrm{ch,top}}$ & Swiss-cheese operad (Vol~II) \\
$\otimes_{\Eone}$ & $\Eone$-factorization tensor:
  $A \otimes_{\Eone} B =
  \mathrm{colim}_{z_1 < z_2} A(z_1) \otimes B(z_2)$.
  \emph{Not} symmetric; IS strictly associative \\
$\otimes_{\Einf}$ & $\Einf$ tensor (Kronecker product):
  kills quantum group structure \\
$\int_M A$ & Factorization homology of $A$ over manifold $M$ \\
\bottomrule
\end{tabular}
\end{center}

%% =====================================================================
\section{CY geometry: K3 and K3$\times E$}
\label{sec:nc-k3}

\begin{center}
\begin{tabular}{@{}lp{8cm}@{}}
\toprule
\textbf{Symbol} & \textbf{Meaning} \\
\midrule
$\Lambda_{\mathrm{Muk}}$ & Mukai lattice of K3, signature $(4,20)$ \\
$\omega^{ij}$ & Mukai pairing: $\mathrm{diag}(+1,\ldots,-1,\ldots)$ \\
$\Phi_{10}$ & Igusa cusp form of weight 10
  (Siegel modular form on $\Sp_4(\Z)$) \\
$\Delta_5$ & BKM automorphic form of weight 5 \\
$\phi_{k,m}$ & Jacobi form of weight $k$, index $m$;
 discriminant constraint $D \equiv 0$ or $3 \pmod{4}$ at $m = 1$ \\
$\eta(\tau)$ & Dedekind eta function: $q^{1/24}\prod_{n \ge 1}(1 - q^n)$ \\
$H_{\mathrm{Muk}}$ & Mukai-lattice Heisenberg VOA:
  $\Phi_2(K3) = H_{\mathrm{Muk}}$, $c = 24$ \\
$\frakg_{K3}$ & K3 double current algebra
  (24 generators, Mukai-signature Serre relations) \\
\bottomrule
\end{tabular}
\end{center}

%% =====================================================================
\section{Holomorphic Chern--Simons}
\label{sec:nc-hcs}

\begin{center}
\begin{tabular}{@{}llll@{}}
\toprule
\textbf{Dimension} & \textbf{Theory} & \textbf{Output} & \textbf{Status} \\
\midrule
3d hol.\ CS & Costello--Gwilliam & Kac--Moody & Proved \\
5d hol.\ CS & Costello 2013 & Affine Yangian & Proved \\
6d hol.\ theory & Costello--Francis--Gwilliam & Quantum toroidal & Conjectural \\
\bottomrule
\end{tabular}
\end{center}

%% =====================================================================
\section{Claim status tags}
\label{sec:nc-claim-status}

\begin{center}
\begin{tabular}{@{}lll@{}}
\toprule
\textbf{Tag} & \textbf{Macro} & \textbf{Meaning} \\
\midrule
\ClaimStatusProvedHere & \verb|\ClaimStatusProvedHere| & Proved in this volume \\
\ClaimStatusProvedElsewhere & \verb|\ClaimStatusProvedElsewhere| & Proved in Vol~I or II \\
\ClaimStatusConjectured & \verb|\ClaimStatusConjectured| & Conjectured \\
\ClaimStatusConditional & \verb|\ClaimStatusConditional| & Conditional on named dependency \\
\ClaimStatusHeuristic & \verb|\ClaimStatusHeuristic| & Heuristic / physical argument \\
\ClaimStatusOpen & \verb|\ClaimStatusOpen| & Open problem \\
\bottomrule
\end{tabular}
\end{center}

\noindent
\textbf{Default rule}: new Vol~III formal statements use
\verb|\begin{conjecture}| unless the proof is complete and unconditional.
Any proof chain passing through an unconstructed object, in particular
$A_X$ at $d = 3$, must remain conjectural, and conditional statuses
propagate: if Result $B$ depends on Result $A$ which is conditional on
an unconstructed object, then $B$ inherits that conditionality.

%% =====================================================================
\section{Main theorems and conjectures}
\label{sec:nc-theorems}

\begin{center}
\begin{tabular}{@{}llll@{}}
\toprule
\textbf{Label} & \textbf{Name} & \textbf{Status} & \textbf{Statement} \\
\midrule
CY-A & CY-to-chiral functor & $d=2,3$ proved; $d \geq 4$ open &
  $\Phi \colon \CY_d\text{-}\Cat \to \En\text{-}\mathrm{ChirAlg}$ \\
CY-B & $\Etwo$-chiral Koszul duality & Programme & Depends on CY-A (resolved) \\
CY-C & Quantum group realization & Conjectural &
  $C(\frakg,q)$ not constructed; $A_\cC$ exists; always \verb|\begin{conjecture}| \\
CY-D & Modular CY characteristic & Programme &
  $A_\cC$ exists by CY-A; $\kappa_\bullet$ formula conjectural at $d = 3$ \\
\bottomrule
\end{tabular}
\end{center}

%% =====================================================================
\section{Deformation parameter \texorpdfstring{$\hbar$}{hbar}}
\label{sec:nc-hbar}

A single $\hbar$ convention is used per file.  The
canonical choice in Vol~III:
\[
  \hbar = \log q, \qquad q = e^\hbar.
\]
If a second convention is needed (e.g.\ $\hbar_{\mathrm{Planck}}$),
a distinct symbol must be introduced with an explicit bridge
identity.
\textbf{Warning}: $\hbar = \log q$ vs $\hbar = (\log q)/(2\pi i)$
differ by a factor of $2\pi i$, which cascades through all formulas.

%% =====================================================================
\section{Degeneration types}
\label{sec:nc-degen}

Different degeneration limits of CY manifolds produce different
chiral algebras with different $\kappa_\bullet$-values.
Every degeneration-limit claim must name the type:

\begin{center}
\begin{tabular}{@{}ll@{}}
\toprule
\textbf{Type} & \textbf{Defining property} \\
\midrule
Large complex structure (LCS) & Maximal unipotent monodromy (MUM) \\
Conifold & Nodal degeneration, $S^3$ cycle collapses \\
Orbifold & Finite group quotient singularity \\
Tropical & Polyhedral degeneration \\
\bottomrule
\end{tabular}
\end{center}
